\chapter{Xerosis}

\section*{Definition}
Xerosis (xeroderma) is abnormally dry skin characterized by scaling, roughness, pruritus, and diminished moisture content of the stratum corneum. It is extremely common, especially in older adults (prevalence >80\% in those \ensuremath{\geq}70 years), and can represent an intrinsic aging change, a manifestation of atopic diathesis, or a sign of systemic disease.

\section*{Pathophysiology}
Xerosis results from reduced stratum corneum water content (<10\% vs. normal 20--30\%), caused by decreased epidermal barrier lipids (ceramides, cholesterol, free fatty acids), impaired filaggrin proteolysis generating natural moisturizing factor (NMF), and altered aquaporin water channel function. The defective barrier increases transepidermal water loss (TEWL) and allows penetration of irritants and allergens, triggering inflammation and pruritus. Low environmental humidity, cold air, frequent bathing, and harsh soaps accelerate moisture loss.

\section*{Etiology}
\begin{itemize}
  \item \textbf{Environmental:} Low humidity (winter, arid climates, central heating), excessive bathing/showering, hot water, harsh soaps and detergents
  \item \textbf{Intrinsic aging:} Reduced sebaceous and eccrine gland activity, decreased NMF, thinning epidermis
  \item \textbf{Atopic dermatitis:} Xerosis is a hallmark feature of atopic diathesis
  \item \textbf{Systemic disease:} Hypothyroidism (myxedema), chronic renal failure, diabetes mellitus, HIV, Sj\"ogren's syndrome, malnutrition (essential fatty acid, zinc, biotin, or vitamin A deficiency)
  \item \textbf{Drug-induced:} Diuretics (thiazides, loop), retinoids (isotretinoin), statins, anticholinergics, cyclosporine
  \item \textbf{Cutaneous disease:} Ichthyosis vulgaris (filaggrin deficiency), lichen simplex chronicus, asteatotic eczema (eczema craquel\'e)
\end{itemize}

\section*{Indications for Evaluation}
Seek care if:
\begin{itemize}
  \item Severe, generalized dry skin with intense pruritus interfering with sleep or daily life
  \item Fissured, cracked skin with bleeding or signs of secondary infection (impetiginization, cellulitis)
  \item Eczema craquel\'e (crazed porcelain appearance) with superficial fissures
  \item Dry skin associated with systemic symptoms: fatigue, weight gain, cold intolerance, polyuria, polydipsia
\end{itemize}

\section*{Related Symptoms}
\begin{itemize}
  \item Pruritus, scaling, fissures, lichenification, asteatotic eczema, ichthyosis
\end{itemize}
\textbf{Note:} Severe, acquired xerosis in adults warrants evaluation for hypothyroidism or chronic kidney disease with TSH, free T4, and basic metabolic panel.

\section*{Self-Care / Home Management}
\begin{itemize}
  \item Apply fragrance-free moisturizing cream (with ceramides, urea 5--10\%, lactic acid 5--12\%, or ammonium lactate 12\%) immediately after bathing and at least twice daily
  \item Limit bathing to 5--10 minutes with lukewarm water once daily; use mild, pH-balanced, fragrance-free cleansers
  \item Use a humidifier (40--50\% humidity) during dry months; pat skin dry instead of rubbing
\end{itemize}

\section*{Diagnostic Approach}
\begin{itemize}
  \item Clinical diagnosis: Dry, rough, scaling skin with accentuation of skin markings; fine white or gray scale; chapping in severe cases
  \item Evaluate for underlying causes: CBC, CMP, TSH/T4, HbA1c, as suggested by history and review of systems
  \item Skin biopsy rarely needed; histology shows compact orthokeratosis, diminished granular layer, and mild perivascular inflammation
\end{itemize}

\section*{Treatment / Management Overview}
\begin{itemize}
  \item First-line: Intensive moisturization with ceramide- or lactic acid-based creams BID
  \item Urea-containing creams (10--20\%) and ammonium lactate 12\% for severe scaling and ichthyotic skin
  \item Fissured or inflammatory xerosis: Short-term topical corticosteroid (hydrocortisone 2.5\% or triamcinolone 0.1\% BID for 5--7 days) added to moisturizer
  \item For asteatotic eczema: Mid-potency topical steroid (triamcinolone 0.1\% cream) for 1--2 weeks with aggressive emollient therapy
  \item Treat underlying systemic condition: thyroid replacement for hypothyroidism, phosphate binders and emollients for uremic xerosis, glycemic control for diabetes
  \item Avoid antipruritic monotherapy (antihistamines alone are ineffective for xerosis without urticaria)
\end{itemize}

\clearpage
